\chapter{Features}

Features are additional model objects that are not defined through element connectivity. They add concentrated properties or auxiliary behaviour to the model.

\section{Point Mass}

\begin{syntax}
*POINTMASS, NSET=<node-set>
<m>, <Ixx>, <Iyy>, <Izz>, <Kx>, <Ky>, <Kz>, <Krx>, <Kry>, <Krz>
\end{syntax}

\cmd{POINTMASS} assigns concentrated mass, rotational inertia, and optional spring stiffness to the nodes of a node set. It is useful for equipment masses, simplified component masses, sensors, local inertias, or concentrated replacement springs.

The data line contains four groups: scalar mass, three rotational inertias, three translational spring stiffnesses, and three rotational spring stiffnesses. Missing values are interpreted as zero.

Point masses can contribute to mass and inertia. They are therefore important for eigenfrequency analysis, transient analysis, inertial loads, and inertia relief. In a purely static loadcase without inertia relief, mass does not automatically mean weight; gravity or acceleration must be introduced through loads.

\cmd{INERTIALOAD} and \cmd{INERTIARELIEF} use the option \key{CONSIDER_POINT_MASSES} to decide whether point-mass features are included in the mass and inertia balance.
